\section{结论与实施建议}

本文用统一的10分钟规划与执行方程连接四个问题，并将“当时可得信息”作为滚动决策的硬边界。核心方法是“合同层滚动线性规划 + 因果贪心反馈”：规划层基于预测轨迹确定合同与目标SOC，执行层在真实量逐槽到达后完成光伏消纳、充放电和应急记账。

第一，典型日连续LP得到购电 $59482.6982\,\mathrm{kWh}$、费用 $35126.9489$ 元。在6000 kWh循环SOC约束下，储能完成低价充电与高价放电，给出了后续全年模型的物理和量纲基线。

第二，固定价格且不能日内调账时，最近60个已完成日期的原始联合残差被用于提取“联合残差特征提取后的边际分位安全轨迹”。正式模式不聚类、$\alpha=0.8$ 事前冻结，所得334日现金费用为1844.64万元，应急购电49.95万kWh。与固定参数的聚类压缩对照相比，未聚类模式在问题二和问题三上均同时满足覆盖率不降低、Pinball损失不增加、费用不增加和应急量不增加，因而被选为正式主线；这项模型仍是确定性分位轨迹LP，不是完整随机MPC或机会约束。

第三，在0、6、12、18时接收官方光伏预报并调整未执行合同时，问题三费用为1622.24万元，应急购电为10.07万kWh；相对问题二减少222.39万元和39.88万kWh，但该总差异同时包含信息来源、调账权和执行路径变化。固定四个重优化时刻、只屏蔽额外官方版本的信息消融表明，全集相对冻结0时节省38.19万元、应急减少2.84万kWh；6、12、18时信息的Shapley贡献分别为12.62、18.88和6.68万元，其中12时贡献最大。这才是额外官方预报版本在当前实验定义下的边际信息价值。

第四，相同0时信息下，固定价2小时合成更新比6小时更新节省62.84万元，即3.8738\%；变价情形节省62.60万元，即3.7551\%。两者分别对应235.17元/次和234.28元/次的增量更新成本阈值。替代结算下，节省分别缩小为0.91万元和2.55万元，说明频率建议高度依赖合同下调条款，不能把主结算的阈值直接外推。

第五，变动电价下，未来价格只由过去30日同槽中位数和当日已实现偏差估计，真实价格只参与事后结算。无调账与有调账策略费用分别为1892.44万元和1667.05万元，应急购电分别为49.89万和10.09万kWh；有调账总体节省225.39万元并减少39.79万kWh应急购电。代表日中仍存在费用和应急同时变差的日期，说明年度优势不等于逐日占优。

数值与稳健性证据形成了完整闭环：5010个规划LP全部达到\texttt{optimal}，最大LP等式残差为 $2.57\times10^{-11}\,\mathrm{kWh}$，最大真实母线平衡残差为 $4.55\times10^{-13}\,\mathrm{kWh}$，规划与执行均未出现同时充放；54组实验、规划--执行SOC偏差和月内经验低光伏分层进一步揭示了模型的适用边界。这些证据支持题设数据内的可行性、因果性和机制解释，但不构成跨年份概率保证。

实施上，建议先把6小时官方发布节奏部署为影子调度，逐日核对预测误差、结算条款、合同闲置和SOC轨迹。在主结算得到确认、数据接入稳定且单次增量更新成本低于约234元时，再试行2小时合成短临更新；数据缺失或估计收益不足时回退到最近一次官方发布方案。若获得多年份天气、节假日和真实高频预报，应采用滚动样本外检验重新评估 $\alpha$、历史窗口和更新频率，并补充储能退化与配电网约束。
